\section{Introduction}
% \elaine{todo: add cites}
\mingxun{todo: add lower bounds}

Private Information Retrieval (PIR),  
originally formulated by Chor, 
Goldreich, Kushilevitz, and Sudan~\cite{pir00}, studies the following important problem.
Imagine that a server holds a public database denoted
${\DB}\in \{0, 1\}^n$.
A client with small local storage  
wants to query the database, while hiding 
its queries from the server. 
PIR has wide applications in practice. For example, 
it enables private contact discovery~\cite{pirpsi,Signalblog},
privacy-preserving light-weight clients for cryptocurrencies, 
private DNS queries~\cite{shi-pir,piano,obliviousdns}, 
private web search~\cite{tiptoe}, and so on. 

\paragraph{Classical PIR without pre-processing.}
A na\"ive solution for PIR is to have the client linearly 
scan through the entire database for each query.
Unfortunately, this would incur linear bandwidth. 
A series of works spanning over two decades~\cite{pir00,pir02,pir06,pir05,pir01,pir07,pir08,pir09,pir10,pir11,pir12,pir13,pir14,pir15} 
starting with Chor et al.~\cite{pir00} showed
%non-trivial constructions of PIR that avoids the linear bandwidth.
how to construct PIR with non-trivial bandwidth. 
Specifically, in the single-server setting, it is well-known that 
with various cryptographic
assumptions (e.g., $\Phi$-hiding, LWE, Damg\r{a}rd-Jurik, DDH, QR), 
we can achieve $\widetilde{O}_\lambda(1)$ bandwidth 
per query~\cite{pir05,simplepir,spiral,tdh}
where $\widetilde{O}_\lambda(\cdot)$ hides  
polylogarithmic factors and the 
dependence on the security parameter $\lambda$.
% \elaine{is it tilde O}
In the two-server setting, 
Dvir and Gopi~\cite{dg16-pir} 
showed that information-theoretic PIR is possible with 
$n^{O(\sqrt{\log \log n /\log n})}$
bandwidth per query.
All the aforementioned 
works studied PIR in the {\it classical setting without preprocessing}. 
More specifically, the classical setting assumes that the server
stores only the original database, and need  
not store per-client state.
Unfortunately, the classical setting 
suffers from the following inherent limitations. 
\begin{enumerate}[leftmargin=5mm]
\item 
First, Beimel, Ishai, and Malkin~\cite{bim-pir}
proved that any classical PIR scheme without preprocessing 
must suffer from 
linear (in $n$) server computation per query. 
Intuitively, if there is any location that the server does not look
at during a query, then the client cannot be asking for that location.
\item 
Second, in a single-server setting, it is known
that any PIR scheme with non-trivial (i.e., sublinear) bandwidth would imply
oblivious transfer~\cite{pirot}, 
i.e., some form of public-key cryptography is needed. 
\end{enumerate}

\ignore{
\paragraph{Pre-processing PIR.}
Both of the above barriers can be overcome 
under a {\it client-specific} preprocessing model, which was first introduced
by Corrigan-Gibbs and Kogan~\cite{sublinearpir}.
The client-specific pre-processing 
model can be thought of as a subscription model. 
During a one-time subscription phase, the client downloads
a hint from the server that is sublinear in $n$, which will
facilitate the client's PIR queries in the future. 
}

\paragraph{Pre-processing sublinear PIR without public-key cryptography.}
To get around the aforementioned barriers,
earlier works have suggested a pre-processing model.
Specifically, we consider a model with a 
one-time pre-processing phase  
upfront, followed by an {\it unbounded} number of queries. 
The pre-processing model was first
proposed by Beimel, Ishai, and Malkin~\cite{bim-pir} and 
Corrigan-Gibbs and Kogan~\cite{sublinearpir}.

In a pre-processing model, 
we focus on exploring the efficiency
of PIR with {\it sublinear computation}, 
and {\it without the use of public-key cryptography} --- 
to get around the aforementioned
barriers,
In particular, the requirement of sublinear computation is especially important
in practical scenarios where the database is large --- for example, in a private
DNS application, the database can be several hundred Gigabytes. 
Further, %understanding how to construct 
we restrict ourselves to Minicrypt (i.e., allowing PRFs but not using any public-key
cryptography) --- this is not only motivated by theoretical interest,  
%such pre-processing PIR schemes without the use of public-key operations is not
%just out of theoretical interest, but also 
but also the promise of concretely faster  constructions since modern processors
have hardware acceleration for AES operations.


%To avoid the linear computation barrier, 
\ignore{
Beimel, Ishai, and Malkin~\cite{bim-pir}
were the first to consider a pre-processing model.
Since then, two types of preprocessing 
have been considered, {\it global server-side} preprocessing~\cite{}, 
and {\it client-specific}
pre-processing~\cite{}. 
In the former (i.e., global server-side preprocessing), 
the server computes some encoding of the database upfront for all clients,
often expanding its storage size to $n^{1+\epsilon}$ or $\poly(n)$.
The latter (i.e., client-specific pre-processing) can be thought of as a subscription model.
Each client subscribes with the server 
often by downloading a hint whose size
is sublinear in $n$. This hint will facilitate the client's PIR queries 
in the future. 
}

Earlier works showed 
that under a global, server-side preprocessing, one can overcome
the linear computation barrier~\cite{bim-pir,sublinearpir,pir-ec22,pir-ec23,treepir,lazza-pir,shi-pir,weikai-depir,piano}; 
and further,
assuming a client-specific preprocessing, we can overcome 
both of the aforementioned barriers~\cite{piano,lingrenpir}. 
So far, we know the following pre-processing PIR schemes
which enjoy  
sublinear pre-processing  
without the use of public-key cryptography.  
Beimel et al.~\cite{bim-pir} showed 
by leveraging a global server-side pre-processing, 
it is possible to construct an 
information theoretic
2-server PIR with $O(n^{1/3})$ bandwidth, $O(n/\log^2 n)$
computation, while consuming $O(n^2)$ 
server storage. The same work also showed
an incomparable information-theoretic 
scheme with $O(n^{1/2 + \epsilon})$
cost in both computation and bandwidth, while incurring
$n^{1+\epsilon'}$
server storage where $\epsilon'$ is a constant dependent on $\epsilon$.
%The recent ground-breaking work of Lin, Mook, and Wich showed how to  
%achieve 
%polylogarithmic bandwidth and computation per query, while consuming  
%$n^{1+\epsilon}$ server storage, assuming the Ring-LWE assumption. 

Under a client-specific pre-processing
model, the recent work \textsc{Piano}~\cite{piano} and the subsequent work of 
Mughees et al.~\cite{lingrenpir} constructed  
single-server PIR schemes with $\widetilde{O}(\sqrt{n})$
bandwidth and server computation, 
$\widetilde{O}_\lambda(\sqrt{n})$
client computation per query, 
while consuming $\widetilde{O}_\lambda(\sqrt{n})$ 
client storage.
Their one-time preprocessing is done by letting the client streamingly download the whole database and dynamically update the client's local state.
Interestingly, their schemes relied
only on one-way functions (OWFs)  
and do not make use of public-key cryptography.
By contrast, classical (single-server) PIR with non-trivial bandwidth implies
oblivious transfer~\cite{pirot}, 
i.e., we cannot get classical single-server PIR in a blackbox
way from OWF~\cite{IR89}.
The reason why Piano~\cite{piano}
 and Mughees et al.~\cite{lingrenpir}
got away using only OWF 
is because their schemes require that the client make a streaming pass  
over
the entire database during preprocessing;  however,
the cost of this preprocessing
can be amortized over an unbounded number of subsequent queries.

\ignore{
Piano~\cite{piano}
 and Mughees et al.~\cite{lingrenpir}'s results
also separate pre-processing PIR
from classical PIR 
(in the single-server setting) 
from a theoretical perspective.
}

\elaine{TODO: for 1server scheme, 
change the asymptotics to the optimized
version}

\begin{table}[t]
\centering
\caption{{\bf Comparison of single-server and two-server pre-processing PIR schemes (for unbounded queries)}.
Any single-server scheme immediately implies a two-server result with the  
same performance bounds. $n$ is the size of the database and $m$ is the number of clients.
The computation overhead counts
both the client and the server's computation, and here we report the expected computation.
The server space counts only the extra storage needed
on top of storing the original database.}
{\scriptsize
\begin{tabular}{cccccccccc}
\toprule
{\bf Scheme} & {\bf Assumpt.} & {\bf Compute} & {\bf Comm.} & \multicolumn{2}{c}{\bf Space} & 
{\bf \#}
& {\bf Concrete}\\
& & & & client & server & {\bf servers} & {\bf eff}.\\
%\cite{lazza-pir} \\
\midrule
\multicolumn{7}{c}{\bf With public-key cryptography}\\[2pt]
\cite{pir-ec22} & LWE & $\widetilde{O}_\lambda(\sqrt{n})$ & 
$\widetilde{O}_\lambda(\sqrt{n})$ & %$\widetilde{O}_\lambda(\sqrt{n})$ &  
$\widetilde{O}_\lambda(\sqrt{n})$ &  
$\widetilde{O}_\lambda(m \cdot n)^*$ & 1 & \xmark \\
%\midrule
\cite{pir-ec23,lazza-pir} & LWE & $\widetilde{O}_\lambda(\sqrt{n})$ & 
$\widetilde{O}_\lambda(1)$ & %$\widetilde{O}_\lambda(1)$ &  
$\widetilde{O}_\lambda(\sqrt{n})$ &  
$\widetilde{O}_\lambda(m \cdot n)^*$ & 1 & \xmark\\
%\midrule
\cite{weikai-depir} & Ring-LWE & $\poly((\log n)^{1/\epsilon})$  & $\poly((\log n)^{1/\epsilon})$ & 
$0$ & $n^{1+\epsilon}$ & 1 & \xmark\\
\cite{shi-pir} & LWE & $\widetilde{O}_\lambda(\sqrt{n})$ &  $\widetilde{O}_\lambda(1)$ &
$\widetilde{O}_\lambda(\sqrt{n})$ & 
0 & 
2 & \xmark \\
\cite{treepir} & Various & $\widetilde{O}_\lambda(\sqrt{n})$ &  $\widetilde{O}_\lambda(1)$ &
$\widetilde{O}_\lambda(\sqrt{n})$ & 
0 & 
2 & \cmark \\[3pt]
\multirow{2}{*}
{\bf Our work}  
& \multirow{2}{*}{Various} & 
\multirow{2}{*}{$\widetilde{O}_\lambda(\sqrt{n})$} & 
{$\widetilde{O}(\sqrt{n})$ offline}
& 
\multirow{2}{*}{$\widetilde{O}_\lambda(\sqrt{n})$} & 
\multirow{2}{*}{$0$} & \multirow{2}{*}{$1$}
%\multirow{2}{*}{$0$} 
& \multirow{2}{*}{\cmark}
\\[-1.5pt]
& & & 
{$\widetilde{O}_\lambda(1)$ online}
& & & \\
\midrule
\multicolumn{7}{c}{\bf Without public-key cryptography}\\[2pt]
\cite{bim-pir} & None & $O(n/\log^2 n)$ & $O(n^{1/3})$ &  $0$ & $O(n^2)$&  2 & \xmark\\
\cite{bim-pir} 
& None & $O(n^{1/2 + \epsilon})$ & $O(n^{1/2 + \epsilon})$ &  $0$ & $O(n^{1+{\epsilon'}})^{**}$ & 2 & \xmark\\[5pt]
%\midrule 
\cite{sublinearpir} & OWF & 
$\widetilde{O}_\lambda(\sqrt{n})$  & 
$\widetilde{O}(\sqrt{n})$  & 
$\widetilde{O}_\lambda(\sqrt{n})$  & 
$0$  & 2 & \cmark \\[5pt]
\cite{checklist} & OWF & 
$O(n)$  & 
$\widetilde{O}_\lambda(1)$  & 
$\widetilde{O}_\lambda(\sqrt{n})$  & 
$0$  & 2 & \cmark \\[5pt]
\cite{piano,lingrenpir} & OWF & 
%& None &  
$\widetilde{O}_\lambda(\sqrt{n})$  & 
${O}(\sqrt{n})$  & 
$\widetilde{O}_\lambda(\sqrt{n})$  & 
$0$  & 1 & \cmark \\[5pt]
%\midrule
%\multirow{2}{*}
{\bf Our work}  
& OWF & 
${O}_\lambda(\sqrt{n})$  & 
${O}_\lambda(n^{1/4})$  & 
$\widetilde{O}_\lambda(\sqrt{n})$  & 
$0$ & $2$
& \cmark
\\[3pt]
%\midrule
\multirow{2}{*}
{\bf Our work}  
& \multirow{2}{*}{OWF} & 
\multirow{2}{*}{${O}_\lambda(\sqrt{n})$} & 
{${O}(\sqrt{n})$ offline}
& 
\multirow{2}{*}{$\widetilde{O}_\lambda(\sqrt{n})$} & 
\multirow{2}{*}{$0$} & \multirow{2}{*}{$1$}
%\multirow{2}{*}{$0$} 
& \multirow{2}{*}{\cmark}
\\[-1.5pt]
& & & 
{${O}_\lambda(n^{1/4})$ online}
& & & \\
\ignore{
    %%%%%%%%%%%%55
[3pt]
\multirow{2}{*}
{\bf Our work}  
& \multirow{2}{*}{LWE or DDH} & 
\multirow{2}{*}{$\widetilde{O}(\sqrt{n})$} & 
{$\widetilde{O}(\sqrt{n})$ offline}
& 
\multirow{2}{*}{$\widetilde{O}(\sqrt{n})$} & 
\multirow{2}{*}{$0$} & \multirow{2}{*}{$1$}
\\[-1.5pt]
& & & 
{$\widetilde{O}_\lambda(1)$ online}
& & & \\
%%%%%%%%%%%%%%%%%%
}
\bottomrule
\end{tabular}
}
{\small $*:$ 
In the unbounded query setting, 
some earlier works~\cite{pir-ec23,lazza-pir,pir-ec22} 
require that the next pre-processing is persistently piggybacked on 
the current window of $O(\sqrt{n})$ operations, and the pre-processing
consumes 
$O_\lambda(n)$ server space
per client to evaluate under FHE an $\widetilde{O}(n)$-sized 
circuit containing a sorting network.

\noindent\small $**:$ 
$\epsilon'>0$ depends on $\epsilon$. 
}
\end{table}

% \elaine{why is treepir LWE or DDH but not various assumptions?}
% \ashrujit{non-preprocessing PIR with linear computation polylog communication are known from only LWE (Spiral) or DDH(Dottling et al Crypto 2019). TreePIR and our new scheme needs to use non-preprocessing PIR as a subroutine}

% \elaine{TODO:change bim desc in abstract}

%\paragraph{Our results.}
\subsection{Our Results}

We show new results that improve the %state-of-the-art 
state of our understanding regarding pre-processing PIR. 
In all of our constructions, the server only needs to store
the original database and need not store any per-client state. 


% \elaine{TODO: fix O lambda issue}

\paragraph{Main result 1.}
First, we construct a two-server pre-processing  
PIR scheme with asymptotically better bandwidth than prior work, relying only on the  
existence of PRFs (which is equivalent to the existence of one-way functions). 
Our result is stated in the following theorem. 

\begin{theorem}[Two-server pre-processing PIR with improved bandwidth]
Assume the existence of one-way functions. 
There exists a two-server pre-processing PIR scheme 
with ${O}_\lambda(n^{1/4})$
bandwidth and ${O}_\lambda(n^{1/2})$
computation per query, while incurring
$\widetilde{O}_\lambda(n^{1/2})$ 
client storage. 
\label{thm:intro-2server}
\end{theorem}
In comparison with the prior work of Beimel et al.~\cite{bim-pir}, 
our \Cref{thm:intro-2server}
achieves significant asymptotic improvements 
in both bandwidth, computation, and server-side storage.
On the other hand, we need to assume one-way functions
whereas Beimel et al.~\cite{bim-pir}'s schemes are information theoretic; further,
we additionally require $\widetilde{O}(\sqrt{n})$ 
space on each client.
However, our construction that gives \Cref{thm:intro-2server} 
is simple and concretely efficient, which is another advantage 
over Beimel et al. %~\cite{}.

\ignore{
If we allow public-key cryptography, 
the TreePIR scheme by Lazzaretti
and Papamanthou achieves 
can further improve the 
bandwidth to $\widetilde{O}_\lambda(1)$ where 
we use $\widetilde{O}_\lambda$ to hide polylogarithmic terms
and dependence on the security parameter $\lambda$.
% \elaine{TODO: finish here.}
}
%However, their computation and client storage costs
%are higher than our scheme because of the dependence 


\paragraph{Main result 2.}
Second, we construct a new pre-processing PIR scheme in the single-server
setting that improves the {\it online} bandwidth in comparison with the
state-of-the-art.  
In this theorem, we differentiate between online bandwidth and offline bandwidth.
The online bandwidth is the bandwidth necessary for the client
to obtain an answer to its query, so it matters to the response time of the client. 
The offline bandwidth is the cost of background maintenance work
amortized to each query, and is not on the critical path of the client's response time. 

\begin{theorem}[Single-server preprocessing PIR with improved online bandwidth]
Assume the existence of one-way functions. 
There exists a single-server pre-processing PIR scheme 
with ${O}_\lambda(n^{1/4})$
{\it online} bandwidth, 
${O}(n^{1/2})$
{\it offline} bandwidth, 
${O}_\lambda(n^{1/2})$ server computation and  
$\widetilde{O}_\lambda(n^{1/2})$ client computation per query, 
while incurring
$\widetilde{O}(n^{1/2})$ 
client storage. 
\label{thm:intro-1server}
\end{theorem}

In comparison with the state-of-the-art scheme \textsc{Piano},
\Cref{thm:intro-1server}
improves the online bandwidth cost
from $\widetilde{O}(\sqrt{n})$ to 
$\widetilde{O}_\lambda(n^{1/4})$, while keeping all other costs
the same. 
Moreover, recall that 
%proven in 
earlier works~\cite{sublinearpir,pir-ec22}, 
proved 
the time-space product lower bound, showing
that the product of the  
client space and the online server time
has to be at least linear in $n$. 
In this sense, 
\Cref{thm:intro-1server}
is tight 
(upto polylogarithmic factors) in terms of this time-space product.

Similar to Piano~\cite{piano} and   
Mughees et al.~\cite{lingrenpir}, 
our 1-server result 
adopts the same model where the client is allowed
to make a streaming pass over 
the database during preprocessing (while consuming
small client space). Otherwise, we would encounter the well-known OT barrier~\cite{pirot}.

We evaluate the concrete performance 
of our 1-server scheme in \Cref{sec:evaluation}.



%\paragraph{Comparison with schemes using public-key cryptography.}



\paragraph{Additional results.}
While our main results focus on constructions in Minicrypt, 
if we are willing to assume classical PIR with $\widetilde{O}_\lambda(1)$ 
bandwidth
(which is known from various assumptions such as 
LWE, $\Phi$-hiding, Damg\r{a}rd-Jurik, DDH, QR)~\cite{pir05,simplepir,spiral,tdh}, 
our techniques would 
then give rise to a concretely efficient
single-server PIR scheme %assuming LWE, 
with 
$\widetilde{O}_\lambda(1)$
{\it online} bandwidth, 
$\widetilde{O}(\sqrt{n})$
offline bandwidth and computation per query, 
consuming $\widetilde{O}_\lambda(\sqrt{n})$
client storage. 
In comparison, although the earlier works
by Zhou et al.~\cite{pir-ec23} and  
 Lazzaretti and Papamanthou~\cite{lazza-pir}
claim to achieve polylogarithmic (online and offline) bandwidth, 
their schemes suffer from a significant drawback, that is, 
the server would have to persistently
store at least $n$ amount of state per client! 
Specifically, Zhou et al.~\cite{pir-ec23} and
 Lazzaretti and Papamanthou~\cite{lazza-pir}
require the pre-processing
phase of the next epoch be piggybacked on 
the queries of the current epoch; however, 
their pre-processing phase requires 
that the server allocate at least $n$ amount of space  
per client, to perform homomorphic evaluation 
of a circuit which is super-linear in size. 
So far, in the unbounded query setting, 
it is not known how to get polylogarithmic overall bandwidth 
(including offline and online)
per query 
{\it under any assumption}, 
assuming that the server stores only the original database.
We state this additional result in the following theorem. 

\begin{theorem}
%[Single-server pre-processing PIR with 
%$\widetilde{O}_\lambda(1)$ 
%online communication]
Assume %\elaine{FILL in various assumptions}.
the existence of a 
classical single-server PIR scheme (i.e., without pre-processing) 
that enjoys $\widetilde{O}_\lambda(1)$ bandwidth per query.
% \elaine{todo: double check}
Then, there exists a single-server pre-processing PIR
scheme with 
$\widetilde{O}_\lambda(1)$ online bandwidth, 
$\widetilde{O}_\lambda(\sqrt{n})$ 
computation, 
$\widetilde{O}(\sqrt{n})$ 
offline bandwidth, 
and requiring 
$\widetilde{O}_\lambda(\sqrt{n})$ 
client storage. 
\label{thm:intro-pkc}
\end{theorem}

% \elaine{TODO: double check the tildes}

 
\subsection{Technical Highlights}
The earlier work of Shi et al.~\cite{shi-pir}
showed that assuming the existence of 
a privately puncturable 
PRF~\cite{privpuncprf00,privpuncprf01,privpuncprf02,privpuncprf03}, 
one can 
construct an efficient 2-server pre-processing PIR scheme
with 
$\widetilde{O}_\lambda(\sqrt{n})$
computation per query and  
requiring 
$\widetilde{O}_\lambda(\sqrt{n})$
client storage. 
Further, the communication per query
is only polylogarithmically larger than
 the size of a punctured key, which can be as small as
$\widetilde{O}_\lambda(1)$ using 
known constructions~\cite{privpuncprf00,privpuncprf01,privpuncprf02,privpuncprf03}. 
Unfortunately, the only known 
techniques 
for constructing a privately puncturable 
PRF~\cite{privpuncprf00,privpuncprf01,privpuncprf02,privpuncprf03}
requires two 
layers of fully homomorphic encryption, and it is not known
whether privately puncturable
PRFs can be built from only one-way functions.  
The elegant TreePIR work of 
Lazzaretti
and Papamanthou~\cite{treepir}
showed how to replace the privately puncturable
PRF with a weaker primitive called a ``weak 
privately puncturable PRF''.
Unfortunately, their approach relies on recursing
on a classical PIR scheme for a $\sqrt{n}$-sized
database, and because this database is {\it dynamically constructed} 
during the scheme, it is not possible 
to pre-process it.
Therefore, Lazzaretti
and Papamanthou~\cite{treepir}'s techniques
fundamentally also require 
public-key cryptography.

\paragraph{Privately programmable pseudorandom set with list decoding.}
Our main contribution is to come up with a new abstraction
called a {\it Privately Programmable 
Pseudorandom Set with List Decoding} (\ppps). 
Given a \ppps key $\sk$, we can 
expand the key $\sk$ to a pseudorandom set 
denoted ${\sf Set}(\sk)$ of size $\sqrt{n}$.
Further, deciding whether any element in $\{0, 1, \ldots, n-1\}$ is in the 
set takes only constant time. 
Importantly, we can call a ${\sf Program}$ 
algorithm to program $\sk$ such that the new  
set  is almost the same as the original ${\sf Set}(\sk)$, except
that the   
one element in the set 
is now changed to another specified element.
%element of choice  
%input into the ${\sf Program}$ algorithm. 
The programmed key 
does not leak information about 
which element is programmed.

Our notion of \ppps is otherwise very similar
to the earlier work of Zhou et al.~\cite{pir-ec23}, except
that we make a relaxation on the correctness when  
decoding a 
programmed key --- this relaxation 
is the crucial
reason why we can construct it from only one-way functions,
whereas Zhou et al.~\cite{pir-ec23}'s construction relies
on LWE.
More specifically,  
we do not require that one can correctly recover
the programmed set given a programmed key $\sk'$.
Instead, we allow {\it list-decoding}, that is, given
a programmed key $\sk'$, decoding outputs a list of candidate
sets, among which one must be the true programmed set.
Moreover, the list-decoding of our PPPS construction is structured, allowing succinct representation and efficient computation.

Using only one-way functions, we construct 
a \ppps scheme with list decoding 
for a pseudorandom set of size $\sqrt{n}$, where
the programmed key 
has size $O_\lambda(n^{1/4})$.


Using such a \ppps scheme, 
we show how to get a two-server
scheme 
with 
${O}_\lambda(n^{1/4})$ communication and  
$\widetilde{O}_\lambda(n^{1/2})$ computation  
per query, using only $\widetilde{O}_\lambda(n^{1/2})$ client space (\Cref{thm:intro-2server}).
Unlike TreePIR~\cite{treepir}, our scheme need not recurse
on a classical PIR scheme, and thus we do not need public-key operations. 

% \elaine{no tilde in table for 2-server, and maybe 1-server too?}


\paragraph{A new broken hint technique.}
\mingxun{transition more smoothly?}
To get our single-server scheme (\Cref{thm:intro-1server}), we encounter
some further challenges. 
In particular, 
it would have been easy to make the scheme work
if our \ppps scheme supported programming a key twice
at two points.
Specifically, in our construction, when the client
consumes a pseudorandom set (represented by $\sk$)
in the hint table
containing the current query $x$, 
it needs to replace the replaced entry with another
randomly sampled \ppps key $\sk$ 
subject to containing the query $x$.
One way to achieve this is to fetch an unconsumed key  
from a backup table, and program the key 
to contain $x$.
However, later, when the client consumes this already-programmed
key in another query $y$, it needs to program  
the point $y$ to 
 some other random point in order not to leak the query $y$.

Unfortunately, our \ppps construction does not support programming twice.
Interestingly, 
earlier works~\cite{pir-ec23,lazza-pir}
also encountered a similar challenge of needing
to program a key twice, but there it was resolved
using different techniques that relied on the LWE 
assumption, which would not work in our setting.

The way we resolve the problem is to introduce
a new technique of allowing {\it broken} entries
in the hint table.
Basically, if the client consumes some \ppps key $\sk$ in the hint table,
it simply replaces the consumed $\sk$ with  
a new entry sampled according to the desired 
distribution (required for privacy). 
However, %since this entry was not prepared in the 
since the client did not perform any preparation work
during the pre-processing phase
for this new entry, 
consuming this entry 
later in a new query 
would result in an incorrect answer, i.e., the replaced entry is {\it broken}.
Fortunately, we can amplify correctness 
through repetition. We defer the details to the subsequent 
technical sections.


\paragraph{Other applications of the broken hint technique.}
The broken hint technique can also lead to other interesting applications.
For example, recall that TreePIR is a 
2-server pre-processing scheme~\cite{treepir}.
With our new broken hint technique, we can convert
TreePIR to a single-server scheme which enjoys the efficiency stated
in \Cref{thm:intro-pkc}.



\paragraph{Further improvements.}
The approach of using broken entries 
introduces a super-logarithmic blowup in the 
bandwidth and computation costs, due the repetition needed
for correctness amplification.
\ifnum\eprintversion=0
In~\Cref{sec:optimized} of the online full version,
\else
In~\Cref{sec:optimized},
\fi
we suggest an improved scheme
that avoids this 
super-logarithmic blowup and gets us the tighter
bounds stated in \Cref{thm:intro-1server}, but 
the resulting scheme is somewhat more complex to describe.

% \elaine{TODO: continue}

%\elaine{TODO: comment on optimality after thm}




\ignore{
\paragraph{Paper Organization.} This paper is organized as follows: in Section~\labelcref{sec:defs} we review formal definitions of single and two-server PIR schemes. In Section~\labelcref{sec:ppps}, we introduce our new abstraction of private programmable pseudorandom sets with list decoding. In Section~\labelcref{sec:two-server}, we give our new two-server preprocessing PIR construction using the aforementioned abstraction that has $O_\lambda(n^{1/4})$ online communication and $O_\lambda(n^{1/2})$ online computation. In Section~\labelcref{sec:one-server} we give our new single-server preprocessing PIR scheme with $\widetilde{O}_\lambda(n^{1/4})$ online communication and $\widetilde{O}_\lambda(n^{1/2})$ online computation. In Section~\labelcref{sec:polylog}, we give a single-server PIR scheme that uses public key cryptography and has polylogarithmic online communication. In Appendix~\labelcref{sec:optimized}, we give an optimized version of the scheme in Section~\labelcref{sec:one-server}, where we open up the black-box of our abstraction to achieve the optimizations.
}

% \subsection{Additional Related Work}

